\documentclass[11pt]{article}
\usepackage[T1]{fontenc}
\usepackage[utf8]{inputenc}
\usepackage{lmodern}
\usepackage[margin=1in]{geometry}
\usepackage{amsmath,amssymb,amsthm,mathtools}
\usepackage{booktabs,array}
\usepackage{hyperref}
\hypersetup{colorlinks=true,linkcolor=blue,urlcolor=blue,citecolor=blue}
\newtheorem{definition}{Definition}
\newtheorem{assumption}{Assumption}
\newtheorem{proposition}{Proposition}
\newtheorem{limitation}{Limitation}
\newcommand{\R}{\mathbb R}
\newcommand{\dd}{\mathrm d}
\newcommand{\Range}{\operatorname{Range}}
\newcommand{\Ker}{\operatorname{Ker}}
\newcommand{\code}[1]{\texttt{\detokenize{#1}}}
\title{Physical Relative Periodic Orbits and Local Branching\\in the Quadruped v3 Hybrid Model}
\author{SLIP Model Zoo: research methods and conditional results}
\date{7 October 2026}
\begin{document}
\maketitle
\begin{abstract}
We define the autonomous compressive-contact quadruped model, its physical
mode manifolds, marked translation-reduced return maps, and the hypotheses
needed for local branching. We derive fixed-foot charts, finite-inertia
mechanical energy conservation, the separate infinite-inertia power identity,
and an exact obstruction to importing delayed scheduled vertical histories.
The propositions are conditional local statements. They establish neither
the requested gait network nor global coverage. Numerical evidence and
remaining application gates are recorded separately in the campaign ledger.
\end{abstract}

\section{Model identity, state, and complete execution}
The declared model identifier is
\code{v3-autonomous-first-directed-root-compressive-stance}. The baseline
parameter vector, in the schema's order, is
\[
p=(k_{l,b},k_{l,f},k_{s,b},k_{s,f},l_{0,b},l_{0,f},
\mathrm{rsla}_b,\mathrm{rsla}_f,j,l_{\mathrm{COM}})^T
=(10,10,20,20,1,1,0,0,2,0.5)^T.
\]
The state is
\[
x=(x_C,\dot x_C,y,\dot y,\phi,\dot\phi,
\alpha_{BL},\dot\alpha_{BL},\alpha_{BR},\dot\alpha_{BR},
\alpha_{FL},\dot\alpha_{FL},\alpha_{FR},\dot\alpha_{FR})^T,
\]
and the separate mode is $q=(q_{BL},q_{BR},q_{FL},q_{FR})^T\in\{0,1\}^4$.
Lengths, body mass, and gravitational acceleration are normalized as in
\code{Quadruped_EOM_v3.md}. All indexes are obtained from
\code{QuadrupedSchema_v3}; a legacy vector is never inserted directly.

\begin{definition}[Autonomous hybrid model]
For fixed $p$, the system consists of modes $Q$, physical domains $D_q$,
tangent vector fields $f_q$, directed guards $G_e$, reset maps $R_e$, event
eligibility, and a simultaneous-event rule. A swing leg is eligible for the
first descending zero of
\[
g_i(x,p)=y+s_i\sin\phi-l_{0,i}\cos(\phi+\alpha_i),
\quad s_i=\begin{cases}-l_{\mathrm{COM}}&i=BL,BR,\\
1-l_{\mathrm{COM}}&i=FL,FR.
\end{cases}
\]
A stance leg is eligible for its first ascending zero. Each reset leaves
body positions and velocities unchanged and projects only the affected leg
rate; its mode bit changes. Distinct-leg simultaneous events form an atomic
batch. Accepted domains require positive stance length, compressive stance,
swing nonpenetration, torso clearance, and the stance-rate constraints below.
\end{definition}

The sufficient state is $(x,q)$; no gait-specific schedule or hidden contact
history is used. A complete execution contains both event-side states,
right-continuous modes, ordered event clusters, and any explicitly declared
auxiliary state. At a section/contact coincidence, the physical batch is
processed once before recording the section. Runge--Kutta stages beyond a
guard used for localization are not accepted hybrid states. The specified
local domains exclude singular geometry and finite-time event accumulation.
Noninvertible resets require local forward semiflows only; no globally
invertible hybrid flow is assumed.

\section{Physical charts and energy}
Set $H_i=y+s_i\sin\phi$, $\theta_i=\phi+\alpha_i$, and
$L_i=H_i/\cos\theta_i$ for stance. Its inferred horizontal foot anchor is
\[
\xi_i=x_C+s_i\cos\phi+H_i\tan\theta_i.
\]
\begin{proposition}[Fixed-foot manifold and local retraction]
On $H_i>0$, $\cos\theta_i>0$, fixed-foot kinematics uniquely determines
\begin{equation}\label{eq:rate}
r_i=-\dot\phi-
\frac{\dot x_C-s_i\sin\phi\dot\phi+
(\dot y+s_i\cos\phi\dot\phi)\tan\theta_i}{H_i\sec^2\theta_i}.
\end{equation}
The $m$ constraints $C_i=\dot\alpha_i-r_i=0$ for stance legs have rank $m$.
The physical mode manifold has dimension $14-m$. Fixing translation and a
regular apex gives dimension $12-m$. Substituting $r_i$ for stance-rate
coordinates is an idempotent local retraction.
\end{proposition}
\begin{proof}
Differentiation gives
\[
\dot\xi_i=\dot x_C-s_i\sin\phi\dot\phi+
(\dot y+s_i\cos\phi\dot\phi)\tan\theta_i+
H_i\sec^2\theta_i(\dot\phi+\dot\alpha_i).
\]
The nonzero last coefficient gives \eqref{eq:rate}. Each constraint has unit
derivative in its own stored rate and no other stored leg-rate dependence,
so the rows are independent. Body coordinates, angles, and swing rates are
independent chart coordinates. The stance acceleration differentiates $r_i$
along $f_q$ on $C=0$, giving $DC_i f_q=0$. Translation and $\dot y=0$ remove
two further independent coordinates. The anchor is reconstructible and
$x_C\mapsto x_C+\ell$ implies $\xi_i\mapsto\xi_i+\ell$.
\end{proof}
On the downward-leg chart an anchor coordinate also recovers
$\alpha_i=\arctan((\xi_i-x_C-s_i\cos\phi)/H_i)-\phi$ locally. The rate chart
avoids redundant anchor state. It fails at zero hip height/leg length or a
horizontal stance leg. A different inverse-tangent branch requires explicit
chart transport. Stored stance rates off $C=0$ are not independent physical
perturbations. In particular, projecting only $\dot y$ to zero does not
produce a physical grounded-apex Floquet stencil.

\begin{proposition}[Mechanical energy and the infinite-inertia distinction]
For finite $j>0$, on the physical manifold,
\begin{equation}\label{eq:energy}
E_q(x)=\frac12(\dot x_C^2+\dot y^2+j\dot\phi^2)+y+
\sum_{q_i=1}\frac12 k_{l,i}(l_{0,i}-L_i)^2
\end{equation}
is conserved along regular modes and zero-compression resets. With $j=\infty$
and $\ddot\phi=0$, the energy obtained by omitting the pitch kinetic term
instead satisfies $\dot E_{\rm red}=-\tau\dot\phi$. It is conserved on the
invariant subspace $\dot\phi=0$, but generally not for a spinning torso.
\end{proposition}
\begin{proof}
For $d_i=(-\sin\theta_i,\cos\theta_i)^T$ and stance force
$F_i=\lambda_i d_i$, fixed-foot geometry gives
$\dot L_i=d_i^Tv_{H,i}$. Spring power is $-\lambda_i\dot L_i=-F_i^Tv_{H,i}$.
Kinetic-plus-gravity power is
$\sum_i F_i^Tv_C+\tau\dot\phi=\sum_i F_i^Tv_{H,i}$, since
$\tau=\sum_i s_i\lambda_i\cos\alpha_i$. These cancel. At TD/LO,
$L_i=l_{0,i}$, so entering/leaving potential is zero; a massless rate reset
does not change body kinetic energy. Omitting pitch kinetic energy removes
the $\tau\dot\phi$ cancellation and yields the stated infinite-$j$ identity.
\end{proof}
Swing oscillators have no finite energy reaction in these infinitesimal-leg
body equations. Adding their formal oscillator energy would create an
incorrect conserved quantity. The expression $\infty\cdot0$ is never
evaluated. If $\phi$ is unwrapped, relative periodic closure at infinite $j$
requires $\dot\phi T=0$, hence $\dot\phi=0$. A winding quotient modulo
$2\pi$ would require a distinct declaration and conservation analysis.

\paragraph{Implementation evidence.}
\code{QuadrupedPhysicalChart_v3} exposes explicit lift, retraction,
coordinate extraction, tangent-basis step-halving diagnostics, and flow
tangency checks. \code{QuadrupedEnergy_v3} returns energy, its analytic
gradient and its flow-power identity. Six MATLAB tests in
\code{TestPhysicalChartEnergy_v3} passed: mode dimensions for all 16 modes,
anchor/rate tangency, tangent actions, finite-$j$ conservation, reset energy,
and the separate infinite-$j$ identity. These are local contracts, not branch
or cluster-regularity certificates.

\section{Relative cycles, minimality, and section charts}
\begin{definition}[Relative periodic execution]
Horizontal translation is $\rho_\ell(x,q)=(x+\ell e_{x_C},q)$ and translates
inferred anchors by $\ell$. A complete execution is relative periodic when
\[
z(t+T)=\rho_\ell z(t),\qquad T>0,
\]
including mode, event sequence, event-side ownership, and auxiliary-state
closure. Its mean speed is $\ell/T$, not generally $\dot x_C(0)$.
\end{definition}
The primitive relative period is the least positive return after translation
reduction. A shorter return with leg relabeling is a spatiotemporal symmetry
unless an additional quotient is declared. A repeated traversal is not a
new primitive orbit; identical geometric traces with different contact
histories need not be equivalent. PIP requires zero drift and synchronous
vertical full-state/contact conditions, not one zero-velocity sample.

The downward-apex geometry is
\[
h(x)=\dot y=0,\qquad Dh f_q=\ddot y<0.
\]
It does not require flight. A BL-relative mark is the last qualifying
downward apex strictly before a specified BL-TD occurrence, with no
intervening BL-TD. The chart records occurrence index, mode, orientation,
cluster ownership, selected apex and return multiplicity. Its identity must
persist for nearby nonperiodic inputs; selection never minimizes periodicity
error. A fixed event-phase shooting gauge is not a state-defined section.
Multiple TD occurrences need explicit marking and chart-overlap transport.

\begin{assumption}[Regular open itinerary neighborhood]\label{ass:regular}
An open neighborhood in independent physical section coordinates preserves
the finite eligible event itinerary, marked return and mode closure. Flows,
domains, directed guards, resets and chart maps are $C^r$, $r\ge1$. Every
event and the terminal section are transverse, geometry is regular, no
earlier enabled guard intrudes, and finite execution holds. At clusters a
compatible regular extension must be separately established on a neighborhood
containing the split-event directions.
\end{assumption}

\begin{proposition}[Return-map regularity]
Under Assumption~\ref{ass:regular} with isolated events, the locally defined
marked return map is $C^r$.
\end{proposition}
\begin{proof}
For a transverse event $g(\varphi_q(t,x))=0$, the implicit-function theorem
gives a local $C^r$ event time with
\[
D\tau(x)=-\frac{Dg\,D_x\varphi_q}{Dg\,f_q}.
\]
Compose variable-time flow, reset and next mode a finite number of times;
apply the same argument to the terminal section time. Compose the physical
lift, translation reduction and output chart extraction. These operations
are $C^r$ on the stated open domain, so their composition is $C^r$.
No inverse reset is used.
\end{proof}
At simultaneous events a continuous finite smooth-selection representation
can instead yield an order-cone-dependent piecewise derivative. Zeroth-order
continuity is needed before claiming a Bouligand derivative. Burden et al.'s
event-selected theorem assumes uniformly transverse event functions and
smooth extensions on an open common vector-field domain
\cite{burden}. Noninvertible resets between different physical manifolds
require an embedding or a direct extension argument; deterministic event
priority does not supply one. Smoothness inside a synchrony subspace does
not establish regularity in daughter directions leaving it.

\begin{proposition}[Fixed points and primitive relative executions]
Assume a single-valued autonomous execution, translation equivariance,
Assumption~\ref{ass:regular}, a consistent mark, and a cycle domain selecting
the specified primitive return. Fixed points of the translation-reduced
marked map correspond to primitive relative periodic executions in that
domain.
\end{proposition}
\begin{proof}
A fixed point closes independent chart coordinates. The physical lift closes
dependent rates; the section/gauge and the policy close modes, events and
auxiliary state. The unreduced terminal state differs only by a translation
$\ell$. Autonomy and equivariance concatenate forward copies translated by
$\ell$, giving relative periodicity. The cycle-domain minimality assumption
excludes shorter complete returns. Conversely, a primitive relative periodic
execution meeting this mark returns after its declared cycle; subtracting
drift and extracting the same chart gives a fixed point.
\end{proof}
\begin{limitation}
One TD/LO per leg and mode closure alone do not prove minimality.
\code{EventCycleReturnPolicy_v3} gives a complete-event contract, with an
independent primitive-period check still required. An iterated map includes
multiple covers; squaring it turns a $-1$ multiplier into $+1$. Neither a
double cover nor a shared geometric grazing limit certifies ordinary
symmetry-breaking ancestry.
\end{limitation}

\section{Fixed-parameter symmetries and isotropy}
\begin{proposition}[Admissible solution transport and chart conjugacy]
Independent BL/BR and FL/FR exchanges preserve the quadruped system for every
admissible fixed $p$, generating $C_2\times C_2$. Horizontal translation is
also a symmetry. A valid symmetry maps admissible complete solutions to
solutions. On a symmetry-invariant marked chart and policy, the induced
action $A_g$ commutes with the return map. Otherwise, a regular phase
transport $J_g$ to the corresponding mark gives the appropriate chart
conjugacy $P_g=J_g P J_g^{-1}$ on its overlap domain.
\end{proposition}
\begin{proof}
Paired legs have identical offsets/parameters, so exchange preserves each
flow equation, physical inequality, guard, eligible transition and affected
rate projection. Apply the symmetry to each ODE segment. Equivariance gives
the transformed segment; compatible guards/eligibility and resets give the
transformed event/post-state. Induction transports the complete execution.
If the mark is invariant, uniqueness gives $P A_g=A_g P$. Otherwise express
the same execution in the transported chart, obtaining the conjugacy. A
coherent single-chart group action additionally requires consistent marked
occurrences and phase composition.
\end{proof}
S4 is generally not a symmetry because front/back hip geometry differs.
When $l_{\rm COM}=1/2$, front/back stiffnesses and lengths agree and
$\mathrm{rsla}_b=-\mathrm{rsla}_f$, combined horizontal reflection and
fore/hind exchange is a symmetry: $x_C,\dot x_C,\phi,\dot\phi,\alpha_i,
\dot\alpha_i$ change sign while front/back labels exchange. Substitution in
hip geometry, forces, torque, swing laws, guards and projection verifies this
claim. At unequal parameters the same operation may map different models;
it cannot be used to identify fixed-$p$ front/hind daughters. No global time
reverser is proved; noninvertible TD rate projection obstructs an inference
based only on velocity-even ODEs.

\begin{definition}[Spatiotemporal isotropy]
For a translation-reduced complete orbit $\bar z$ of primitive period $T$,
\[
H_{\bar z}=\{(g,\vartheta):g\bar z(t)=\bar z(t+\vartheta T)
\text{ for every }t,\ \vartheta\in\R/\mathbb Z\}.
\]
Pointwise spatial isotropy has $\vartheta=0$; reversing symmetry is separate.
\end{definition}
Equal contact timings do not establish full-state isotropy. Subgroup loss
up to conjugacy is reported only after full-trajectory/mode transport is
checked. Maximal isotropy need not identify one unique PIP root family.

\section{Local branching, multiple critical modes, and missing hypotheses}
\begin{proposition}[A sufficient simple local branching theorem]\label{prop:cr}
Let the centered independent residual satisfy $F(0,\mu)=0$ and be $C^2$
on an open neighborhood of $(0,\mu_*)$. Assume
$L=D_uF(0,\mu_*)$ is Fredholm of index zero, with one-dimensional kernel
$\operatorname{span}\{v\}$ and one-dimensional cokernel, and
\[
D_{\mu u}F(0,\mu_*)v\notin\Range L.
\]
Then a nontrivial local solution curve exists with leading state direction
$v$. This is a sufficient regularity formulation of the centered
Crandall--Rabinowitz theorem; the original theorem 1.7 specifies weaker
partial-derivative requirements \cite{cr}.
\end{proposition}
\begin{proof}[Proof sketch]
Split $u=av+w$ into kernel and a complement, and the codomain into range
and cokernel. The implicit-function theorem solves the range equation for
$w=w(a,\mu)$. Divide the scalar cokernel equation by $a$ and extend at
$a=0$. Transversality gives a nonzero derivative in $\mu$, so another
implicit-function argument determines $\mu(a)$ and the daughter curve.
\end{proof}
A pitchfork additionally needs an involution acting by $a\mapsto-a$ on the
critical sector, compatible residual equivariance, higher regularity and
nondegenerate reduced coefficients:
\[
0=a\{c(\mu-\mu_*)+d a^2+\cdots\},\qquad cd\ne0.
\]
Then the two reflected daughters have quadratic parameter-amplitude
scaling. A $+1$ multiplier verifies neither those coefficients nor the
smoothness/transversality assumptions. At an event-order boundary the $C^2$
extension must include adjacent admissible daughter itineraries. Otherwise
the candidate requires a justified nonsmooth argument and remains unresolved.

For a multiple kernel $K$, range reduction leaves
$\Psi(a,\mu)=0$ on $K$ with its induced symmetry representation. Decompose
$K$ into sectors and evaluate the reduced nonlinear coefficients. A sector
restriction is valid only if the full residual/map preserves it and the
restricted kernel and cokernel meet Proposition~\ref{prop:cr}. The protected
P1 pronking reference reports three tangent-free bounding/front-spread/
hind-spread critical modes. Its six local numerical rays do not reduce the
entire critical space to one simple eigenvector and do not establish an
autonomous v3 network.

\section{Physical linearization and structural neutrality}
For an isolated transverse autonomous event, common-time comparison of
perturbed event executions gives the saltation matrix
\begin{equation}\label{eq:saltation}
S_e=DR_e+\frac{(f^+-DR_e f^-)n_e^T}{n_e^Tf^-},
\qquad n_e=\nabla g_e(x^-).
\end{equation}
The assumptions include differentiable guards/resets, transversality and
stable local event selection \cite{kong}. The reset Jacobian alone omits
event-time sensitivity. Compose flow variational matrices and saltations
in chronological order, apply the terminal section correction
\[
\Pi_\Sigma=I-\frac{f\,n_\Sigma^T}{n_\Sigma^Tf}
\]
once, then apply the actual translation gauge and physical chart extraction.
Initial perturbations enter through the derivative $B$ of the physical
lift. A finite difference of a redetected return map already includes timing;
adding another saltation/terminal timing correction double-counts it.

When regular augmented shooting guards satisfy $H(x,\tau)=0$,
\[
D_x\tau=-H_\tau^{-1}H_x,\qquad
D_xF_{\rm eff}=F_x-F_\tau H_\tau^{-1}H_x.
\]
The same lifting translates a physical critical direction into augmented
solver coordinates. Eigenvalues of the joint shooting residual Jacobian are
not physical Floquet multipliers.

If $E(P(u))=E(u)$ and $DE\ne0$, differentiation at a fixed point gives
\[
DE\,DP=DE.
\]
Energy supplies a structural left unit eigenvector. A differentiable fixed-$p$
family of fixed points supplies $DP\,u_s=u_s$. Varying $p$ adds a $P_p p_s$
term and does not automatically supply a state-only unit right vector.
Phase and translation are removed by section/gauge. A fixed-energy chart
removes one coordinate and a dependent closure equation by rank; blindly
adding an energy equation to redundant periodicity equations is unjustified.
No fixed number of eigenvalues closest to one is discarded.

At a cluster, retain feasible order-cone-dependent derivative selections
only after continuity and piecewise-regularity hypotheses hold. Commuting
instantaneous resets do not prove equal full-flow/time-sensitivity products
or $C^2$ regularity. Individual selection-matrix spectra do not establish
stability under cone transitions; a justified common contraction norm is
a sufficient alternative. Record matrix-action/step/tolerance convergence,
not eigenvalue agreement alone.

\section{Exact vertical histories and a specific incompatibility}
At baseline vertical synchrony,
\[
\ddot y=40(1-y)-1\quad\text{in stance},\qquad
\ddot y=-1\quad\text{in flight}.
\]
Write $\omega=\sqrt{40}$, $y_{\rm eq}=39/40$, $a=1/40$,
$v=\sqrt{2(E-1)}$ and $\theta=\operatorname{atan2}(v/\omega,a)$.
Starting at descending TD, $y=1$, $\dot y=-v$,
\[
y(t)=y_{\rm eq}+a\cos(\omega t)-(v/\omega)\sin(\omega t).
\]
The first ascending $y=1$ exit occurs at
$S_0=(2\pi-2\theta)/\omega$; flight lasts $2v$. A scheduled model retaining
$n$ extra complete stance revolutions has
\begin{equation}\label{eq:vertical}
T_n(E)=\frac{2\pi-2\theta}{\omega}+2v+n\frac{2\pi}{\omega},
\qquad n=0,1,2,\ldots.
\end{equation}
The minimum height is
$y_{\rm eq}-\sqrt{a^2+2(E-1)/40}$, so positive leg length requires
$1<E<20$. For every $n\ge1$, the additional revolution reaches
$y_{\max}=y_{\rm eq}+\sqrt{a^2+2(E-1)/40}>1$, giving tensile stance.
More immediately, the autonomous first-root law has already executed LO
at $S_0$. Delayed histories therefore cannot be imported into this model
without changing eligibility or domain. Each such schedule has one flight
and extra stance oscillations, not a repeated complete ordinary cycle.

As $E\downarrow1$, $T_n\to(n+1)2\pi/\omega$ while the common all-stance
geometric oscillator has minimal period $2\pi/\omega$. Transversality
vanishes and delayed representations become multiple covers. This is a
grazing limit, not a proved regular bifurcation. The obstruction concerns
that scheduled vertical route; it does not prove that no admissible branch
can leave synchrony and reach another B2 orbit elsewhere.

\section{Research meaning and limitations}
The registered root is the ordinary first-LO PIP family. The distinct
claims are specific local connections, a single-root target network,
bounded-domain target-class coverage, and the stronger global conjecture
that every primitive orbit belongs to the component. Only validated edges
of the declared type count. Imported-seed replay is distinct from frozen
parent-only discovery. Different parameter/model identities cannot be joined
to complete a fixed-model graph.

The protected B2 G/E records are alternative apex views, not separate
discoveries. Legacy suffix 2 does not certify flight count. The P2 historical
route begins on delayed PIP and includes suppressed roots/tension; ordinary
PIP-to-delayed PIP is not a verified regular edge. Local propositions cannot
imply global connectivity, unique maximal roots, exhaustive discovery or
dynamic stability. Disconnected components, higher primitive periods,
secondary bifurcations, missing symmetry types and singular itinerary
boundaries remain possible. A failed corrector or budget stop is unresolved,
not a proof of nonexistence. Global claims require additional certification.

\begin{thebibliography}{9}
\bibitem{cr}
M. G. Crandall and P. H. Rabinowitz,
``Bifurcation from Simple Eigenvalues,''
\emph{Journal of Functional Analysis} 8 (1971), 321--340.
\href{https://doi.org/10.1016/0022-1236(71)90015-2}{doi:10.1016/0022-1236(71)90015-2}.
Original theorem 1.7 inspected at
\href{https://jxshix.people.wm.edu/2013-taiwan/crandall-rabinowitz-1971-JFA.pdf}{author-hosted scan}.

\bibitem{kong}
N. J. Kong, J. J. Payne, J. Zhu and A. M. Johnson,
``Saltation Matrices: The Essential Tool for Linearizing Hybrid Dynamical
Systems,'' \emph{Proceedings of the IEEE} 112(6) (2024), 585--608.
\href{https://arxiv.org/html/2306.06862v3}{Author manuscript, arXiv:2306.06862v3}.

\bibitem{burden}
S. A. Burden, S. S. Sastry, D. E. Koditschek and S. Revzen,
``Event-Selected Vector Field Discontinuities Yield Piecewise-Differentiable
Flows.''
\href{https://arxiv.org/html/1407.1775v3}{Author manuscript, arXiv:1407.1775v3}.
Definitions 1--2 and theorem 3 inspected. The theorem is not a blanket result
for arbitrary noninvertible resets or grazing.
\end{thebibliography}
\end{document}
